% notes/08_benign_overfitting_research_programme.tex
%
% Research programme — transductive benign overfitting for the min-norm
% RKHS interpolant on a finite graph.
%
% This document supersedes an earlier draft (notes/08_benign_overfitting_NTK.tex,
% v1) which contained mathematically defective "proof sketches"
% identified by an adversarial review (Gemini-Pro-2.5, pre-submission
% pass). The defects were:
%
%   (D1) Confusion between l_2(V) norm and RKHS norm in the bias
%        bound (Lemma 7 of v1): the inequality
%        ||b_T(y_>r)||_l2 <= ||y_>r||_l2 is FALSE for the min-norm
%        RKHS interpolant, which is only non-expansive in RKHS norm.
%
%   (D2) Category error in the variance bound (Theorem 8 of v1):
%        BLLT Lemma 10 is stated for empirical covariance matrices of
%        i.i.d. Gaussian designs, not for principal submatrices of a
%        fixed Gram K extracted by without-replacement sampling. The
%        correct tool is matrix Chernoff for principal submatrices
%        with leverage-score / incoherence conditions.
%
%   (D3) Conditional-expectation error in the transductive Bernstein
%        step (Lemma 9 + Theorem 8 of v1): the trace argument
%        Tr(K(T^c,T) K_T^{-2} K(T,T^c)) depends on T, so
%        Bardenet-Maillard on the diagonal is invalid for it.
%
% Rather than paper over these defects with a half-corrected v2,
% this version reframes notes/08 as a RESEARCH PROGRAMME:
%   - The bias-variance decomposition (Sections 2-3) is correct and
%     retained verbatim.
%   - The "Theorems" of v1 are demoted to "Conjectures" and stated
%     with explicit dependence on the missing technical ingredients.
%   - A new Section 5 lists the precise mathematical tools needed to
%     prove each conjecture, with pointers to the relevant literature
%     and an honest assessment of difficulty.

\documentclass[11pt,a4paper]{article}

\usepackage[margin=2.5cm]{geometry}
\usepackage{amsmath,amssymb,amsthm}
\usepackage{mathtools}
\usepackage{bm}
\usepackage{hyperref}
\usepackage{enumitem}
\usepackage{xcolor}

\usepackage[authoryear,round]{natbib}
\bibliographystyle{plainnat}

\theoremstyle{plain}
\newtheorem{theorem}{Theorem}
\newtheorem{proposition}[theorem]{Proposition}
\newtheorem{lemma}[theorem]{Lemma}
\newtheorem{corollary}[theorem]{Corollary}
\newtheorem{conjecture}[theorem]{Conjecture}
\theoremstyle{definition}
\newtheorem{definition}[theorem]{Definition}
\newtheorem{assumption}[theorem]{Assumption}
\newtheorem{problem}[theorem]{Open Problem}
\theoremstyle{remark}
\newtheorem{remark}[theorem]{Remark}

\newcommand{\by}{\mathbf{y}}
\newcommand{\bxi}{\boldsymbol{\xi}}
\newcommand{\bbeta}{\boldsymbol{\beta}}
\newcommand{\balpha}{\boldsymbol{\alpha}}
\newcommand{\bphi}{\boldsymbol{\phi}}
\newcommand{\bP}{\mathbf{P}}
\newcommand{\bB}{\mathbf{B}}
\newcommand{\bI}{\mathbf{I}}
\newcommand{\bK}{\mathbf{K}}
\newcommand{\bM}{\mathbf{M}}
\newcommand{\bX}{\mathbf{X}}
\newcommand{\bV}{\mathbf{V}}
\newcommand{\bU}{\mathbf{U}}
\newcommand{\bLambda}{\boldsymbol{\Lambda}}
\newcommand{\bone}{\mathbf{1}}
\newcommand{\bv}{\mathbf{v}}
\newcommand{\bu}{\mathbf{u}}
\newcommand{\bz}{\mathbf{z}}
\newcommand{\bg}{\mathbf{g}}
\newcommand{\bb}{\mathbf{b}}
\newcommand{\E}{\mathbb{E}}
\newcommand{\R}{\mathbb{R}}
\newcommand{\PR}{\mathbb{P}}
\newcommand{\Var}{\mathrm{Var}}
\newcommand{\Cov}{\mathrm{Cov}}
\newcommand{\Rspec}{R^2_{\mathrm{spec}}}
\newcommand{\Scal}{\mathcal{S}}
\newcommand{\Acal}{\mathcal{A}}
\newcommand{\Hcal}{\mathcal{H}}
\newcommand{\Rcal}{\mathcal{R}}
\newcommand{\Ycal}{\mathcal{Y}}
\newcommand{\Vcal}{\mathcal{V}}
\newcommand{\Dcal}{\mathcal{D}}
\newcommand{\KL}{\mathrm{KL}}
\newcommand{\TV}{\mathrm{TV}}
\newcommand{\rank}{\mathrm{rank}}
\newcommand{\tr}{\mathrm{tr}}
\newcommand{\df}{\mathrm{df}}
\newcommand{\reff}{r_{\mathrm{eff}}}
\newcommand{\Reff}{R_{\mathrm{eff}}}
\newcommand{\range}{\mathrm{range}}
\DeclareMathOperator*{\argmin}{arg\,min}

\title{\textbf{Transductive benign overfitting: a research programme}\\
\large{Identifying the technical obstacles to a leave-node-out
analogue of Bartlett--Long--Lugosi--Tsigler}}
\author{Rohan Foss\'e \and Ga\"el Pallares\\
\small{CESI LINEACT (EA 7527), Montpellier, France}}
\date{Working notes, May 2026 --- \emph{research programme, not a
proven theorem}}

\begin{document}

\maketitle

\begin{abstract}
Theorem~4 of \texttt{notes/07\_four\_axes\_extensions.tex} (the
universal spectral bound on the $\Omega$-effective reachable
subspace) becomes \emph{operationally vacuous} for the min-norm
RKHS interpolating learner on a strictly positive-definite kernel
$\bK$, since then $\range(\bK) = \R^N$ and the spectral $R^2$ floor
is~$1$. Adapting the \citet{BartlettLongLugosiTsigler2020}
benign-overfitting analysis to recover a non-trivial leave-node-out
floor is a natural research direction, but \emph{not} a routine
adaptation: BLLT's machinery rests on i.i.d.\ Gaussian designs and
empirical covariance concentration, neither of which applies to the
principal submatrix of a fixed Gram matrix extracted by
without-replacement sampling.

This note formulates the transductive benign-overfitting programme
as a precise sequence of \textbf{conjectures}, each one explicitly
parameterised by the technical ingredient it needs in order to
become a theorem. The bias--variance decomposition of the
held-out risk (\S\ref{sec:decomp}) is correct and unconditional;
the bias bound (\S\ref{sec:bias-conjecture}) requires a correct
RKHS-vs-$\ell_2(V)$ norm conversion that exploits the spectral
structure of $\bK_T^+$; the variance bound (\S\ref{sec:var-conjecture})
requires (i)~a matrix Chernoff inequality for principal submatrices
sampled without replacement and (ii)~an exchangeability-based
concentration argument for trace functionals of $T$-dependent random
matrices. \S\ref{sec:tools} catalogues the missing tools and points
to the relevant literature. \S\ref{sec:application} discusses when
the BLLT regime can be expected to hold on graphs.

An earlier draft (\texttt{notes/08\_benign\_overfitting\_NTK.tex},
v1) contained ``proof sketches'' for the two key Conjectures
that papered over these obstacles by misapplying BLLT~Lemma~10 and
\citet{BardenetMaillard2015}; we deliberately retract those sketches
here and replace them with this explicit research programme. The
intellectual contribution of this note is therefore not new
mathematics, but the precise identification of \emph{what needs to
be proved} and \emph{by which tools}.
\end{abstract}

\tableofcontents

% =========================================================================
\section{Setup}
\label{sec:setup}

We retain the notation of \texttt{notes/01}~\S2 and
\texttt{notes/07}~\S1.

\subsection{The min-norm RKHS interpolant}

\begin{definition}[Min-norm RKHS interpolant]
\label{def:min-norm-interp}
Let $K : V \times V \to \R$ be a symmetric positive-semidefinite
kernel with Gram matrix $\bK \in \R^{N \times N}$ on $V$. For a
training set $T \subset V$ of size $n$ and observations $\bz \in \R^n$
indexed by $T$, the \emph{min-norm RKHS interpolant} is
\begin{equation}
    \hat f(\,\cdot\,; T, \bz)
    \;:=\; \argmin\!\big\{\|f\|_K^2 : f \in \range(\bK),\;\; f|_T = \bz\big\}
    \;=\; \bK(\cdot, T)\,\bK_T^{+}\,\bz,
    \label{eq:min-norm-interp}
\end{equation}
where $\bK_T := \bK(T, T) \in \R^{n \times n}$ and $\bK_T^+$ is its
Moore--Penrose pseudo-inverse. The norm in the optimisation is the
RKHS norm $\|f\|_K^2 := \langle f, \bK^+ f\rangle$, \emph{not} the
Euclidean norm.
\end{definition}

\begin{remark}[Reduction to NTK / SGD-trained NN]
\label{rem:ntk-reduction}
For a sufficiently wide NN with $\Theta$-Lipschitz activation
trained by SGD from small initialisation, the learned predictor
converges to the min-norm interpolant of the empirical NTK
\citep{JacotGabrielHongler2018,LeeXiaoSchoenholzBahriNovakSohlDicksteinPennington2019}.
We work abstractly with any PSD kernel $K$.
\end{remark}

\subsection{Spectral structure}

Let $\bK = \bV\bLambda\bV^\top$ be the eigendecomposition with
$\bLambda = \mathrm{diag}(\hat\lambda_1, \ldots, \hat\lambda_N)$,
$\hat\lambda_1 \ge \cdots \ge \hat\lambda_N \ge 0$, and
$\bV = [\bv_1, \ldots, \bv_N]$ orthonormal. Let
$\Scal_k := \mathrm{span}(\bv_1, \ldots, \bv_k)$ and $\bP_{\Scal_k}$
its orthogonal projector.

\begin{definition}[BLLT effective rank and dimension]
\label{def:eff-rank}
For $k \in \{0, \ldots, N-1\}$:
\begin{align}
    \reff_k(\bK) &:= \tfrac{(\sum_{i>k} \hat\lambda_i)^2}{\sum_{i>k} \hat\lambda_i^2},
    \qquad
    \Reff_k(\bK) := \tfrac{\sum_{i>k} \hat\lambda_i}{\hat\lambda_{k+1}}.
    \label{eq:eff-rank}
\end{align}
\end{definition}

\begin{definition}[Leverage scores]
\label{def:leverage}
The \emph{leverage score} of node $v \in V$ at level $k$ is
$\ell^{(k)}_v := \sum_{j=1}^k (\bv_j)_v^2 \in [0, 1]$. The
\emph{leverage incoherence} of $\bV$ at level $k$ is
$\mu_k(\bV) := \max_{v \in V} \ell^{(k)}_v / (k/N)$. We have
$\mu_k(\bV) \in [1, N/k]$, with $\mu_k = 1$ in the fully
incoherent case (Fourier basis on a regular graph).
\end{definition}

\begin{definition}[Spectral tail of the target]
\label{def:spec-tail-y}
For $r \in \{0, \ldots, N\}$,
$\tau_r(\by) := \|\bP_{\Scal_r^\perp}\by\|^2/\|\by\|^2 \in [0, 1]$.
The target is \emph{$(r, \tau)$-spectrally concentrated} if
$\tau_r(\by) \le \tau$.
\end{definition}

\subsection{Noise model}

\begin{assumption}[Sub-Gaussian noise]
\label{a:noise-bllt}
$\bxi \in \R^N$ has independent zero-mean entries with
$\Var(\xi_v) = \sigma_\xi^2$, $\sigma_\xi$-sub-Gaussian
componentwise, independent of $T \sim \Pi_{\mathrm{LSO}}$.
\end{assumption}

% =========================================================================
\section{Decomposition of the held-out error (unconditional)}
\label{sec:decomp}

The following decomposition is exact and does not depend on any
benign-overfitting hypothesis.

\subsection{Bias-variance split}

By linearity of \eqref{eq:min-norm-interp} in $\bz$, the
interpolant splits as
\begin{equation}
    \hat f^{\mathrm{interp}}_\Acal(T)
    \;=\; \underbrace{\bK(\cdot, T)\bK_T^+ \by|_T}_{=:\, \bb_T(\by)\quad \text{(``bias'' map)}}
    \;+\; \underbrace{\bK(\cdot, T)\bK_T^+ \bxi|_T}_{=:\, \bg_T(\bxi)\quad \text{(``variance'' map)}}.
    \label{eq:bv-split}
\end{equation}
By independence of $\bxi$ and $(T, \by)$, the held-out squared
error decomposes (conditionally on $T$) as
\begin{equation}
    \E_\bxi\!\big[L_{T^c}(\hat f^{\mathrm{interp}}_\Acal(T), \by) \,|\, T\big]
    \;=\; \underbrace{L_{T^c}(\bb_T(\by), \by)}_{=:\, \mathrm{Bias}_{T^c}^2(\by, T)}
        \;+\; \underbrace{\E_\bxi\!\big[L_{T^c}(\bg_T(\bxi), 0) \,|\, T\big]}_{=:\, \mathrm{Var}_{T^c}^2(\sigma_\xi^2, T)}.
    \label{eq:bv-decomp}
\end{equation}

\begin{lemma}[Interpolation $\Rightarrow$ bias supported on $T^c$]
\label{lem:bias-on-tcomplement}
For any $\by \in \range(\bK)$, the bias map satisfies
$\bb_T(\by)|_T = \by|_T$. Hence
$\mathrm{Bias}_{T^c}^2(\by, T) = |T^c|^{-1}\|\bb_T(\by) - \by\|_{T^c}^2$,
and the bias error of the interpolant on $V$ is concentrated on $T^c$.
\end{lemma}

\begin{proof}
We need $\by|_T \in \range(\bK_T)$. If $\by = \bK\balpha$ for some
$\balpha \in \R^N$, then
$\by|_T = \bK(T, V)\balpha = \bK_T \balpha_T + \bK(T, T^c)\balpha_{T^c}$.
Since $\bK$ is PSD, the Schur complement structure of $\bK$
(positivity of the $2\times 2$ block $\binom{\bK_T \bK(T,T^c)}{\bK(T^c,T) \bK_{T^c}}$)
implies $\range(\bK(T, T^c)) \subseteq \range(\bK_T)$, so
$\bK(T, T^c)\balpha_{T^c} \in \range(\bK_T)$, hence
$\by|_T \in \range(\bK_T)$. Then $\bK_T \bK_T^+ \by|_T = \by|_T$
since $\bK_T\bK_T^+$ is the orthogonal projector onto $\range(\bK_T)$.
The first claim of the lemma follows.
For the second: $\bb_T(\by)|_T = \by|_T$ implies the pointwise
error $\bb_T(\by)(v) - y_v = 0$ on $T$, hence the sum
$\sum_{v \in V}(\bb_T(\by)(v) - y_v)^2 = \sum_{v \in T^c}(\bb_T(\by)(v) - y_v)^2$.
Dividing by $|T^c|$ gives $\mathrm{Bias}_{T^c}^2$.
\qed
\end{proof}

\begin{lemma}[Spectral representation of the bias map]
\label{lem:bias-spectral-rep}
Let $\bV_T := \bV(T, :) \in \R^{n \times N}$. Then
\begin{equation}
    \bb_T(\by)
    \;=\; \bV\bLambda\bV_T^\top (\bV_T\bLambda\bV_T^\top)^+ \by|_T.
    \label{eq:bias-spectral}
\end{equation}
\end{lemma}

\begin{proof}
Direct substitution into $\bK(\cdot, T)\bK_T^+ \by|_T$ using
$\bK = \bV\bLambda\bV^\top$, $\bK(\cdot, T) = \bV\bLambda\bV_T^\top$,
$\bK_T = \bV_T\bLambda\bV_T^\top$. \qed
\end{proof}

\subsection{Spectral representation of the variance term}

\begin{lemma}[Closed-form variance trace]
\label{lem:var-spectral-rep}
Under A\ref{a:noise-bllt}, for fixed $T$:
\begin{equation}
    \E_\bxi\!\big[L_{T^c}(\bg_T(\bxi), 0) \,|\, T\big]
    \;=\; \frac{\sigma_\xi^2}{|T^c|} \tr\!\Big(\bK(T^c, T)\bK_T^+\, \bK_T^+\, \bK(T, T^c)\Big).
    \label{eq:var-rep}
\end{equation}
\end{lemma}

\begin{proof}
$\E_\xi[\bg_T(\bxi)\bg_T(\bxi)^\top \mid T]
= \sigma_\xi^2 \bK(T^c, T)\bK_T^+\bK_T^+\bK(T, T^c)$; restrict to
$T^c$, take trace, divide by $|T^c|$. \qed
\end{proof}

\textbf{All material up to this point is unconditional and correct.}
The conjectural part begins here.

% =========================================================================
\section{Conjecture: bias bound}
\label{sec:bias-conjecture}

\subsection{The norm conversion problem}

A natural attempt is to bound the bias on $T^c$ by decomposing
$\by = \by_{\le r} + \by_{>r}$ with $\by_{\le r} := \bP_{\Scal_r}\by$
and arguing:
\begin{enumerate}[label=(\alph*),leftmargin=*]
\item \emph{Head contribution.} $\bb_T(\by_{\le r})$ interpolates
    a signal in $\Scal_r$; by the Pesenson sampling-quality
    argument (\texttt{notes/01}~Lemma~3), the held-out error of
    the interpolant on $\Scal_r$ is bounded by a slack
    depending on $\sigma_\Pi(\Scal_r)$.
\item \emph{Tail contribution.} $\bb_T(\by_{>r})$ is the
    interpolant of a high-frequency signal; we need to bound
    $\|\bb_T(\by_{>r}) - \by_{>r}\|_{T^c}^2/|T^c|$.
\end{enumerate}

\textbf{The defect in v1 of these notes.}  An earlier draft argued
$\|\bb_T(\by_{>r})\|_2^2 \le \|\by_{>r}\|_2^2$ ``by minimum-norm''.
This is \emph{false}: the min-norm RKHS interpolant minimises the
RKHS norm $\|f\|_K^2 = \langle f, \bK^+f\rangle$, \emph{not} the
Euclidean $\ell_2(V)$ norm. The conversion factor is
$\|f\|_2^2 \le \hat\lambda_1 \|f\|_K^2$ in one direction, but the
reverse $\|f\|_K^2 \le \hat\lambda_N^{-1}\|f\|_2^2$ blows up when
$\bK$ has small eigenvalues --- exactly the regime where $\by_{>r}$
lives.

\subsection{What the correct bound looks like}

\begin{conjecture}[Bias bound, transductive analogue of BLLT~\S5.1]
\label{conj:bias}
Suppose $\by$ is $(r, \tau)$-spectrally concentrated, the leverage
incoherence at level $r$ is bounded ($\mu_r(\bV) \le \mu^*$ for some
absolute constant), and the Pesenson sampling quality
$\sigma_\Pi(\Scal_r) \ge 1/2$. Then there exist absolute constants
$C_1, C_2$ such that
\begin{equation}
    \E_T\!\big[\mathrm{Bias}_{T^c}^2(\by, T)\big]
    \;\le\; C_1\,\tau\,\Var(\by)
        \;+\; C_2\,\frac{r\,\mu^*}{n - r}\,(1 - \tau)\,\Var(\by).
    \label{eq:bias-conjecture}
\end{equation}
\end{conjecture}

\begin{remark}[On the conjecture vs.\ the v1 defective ``proof'']
\label{rem:bias-discussion}
Conjecture~\ref{conj:bias} differs from the v1 ``Lemma~7''
statement \eqref{eq:bias-conjecture} in two crucial ways:
\begin{enumerate}[label=(\arabic*),leftmargin=*]
\item The constant $C_1$ in front of $\tau\Var(\by)$ is \emph{not}
    necessarily $1$. The tail contribution
    $\|\bb_T(\by_{>r}) - \by_{>r}\|_{T^c}^2/|T^c|$ is bounded by
    $\tau\Var(\by)$ \emph{only if} the interpolation amplification
    of $\bK_T^+$ on tail frequencies can be controlled. The norm
    conversion $\|\bb_T(\by_{>r})\|_K^2 \le \|\by_{>r}\|_K^2$
    \emph{is} valid (RKHS non-expansion), giving
    $\|\bb_T(\by_{>r})\|_2^2 \le \hat\lambda_1\|\by_{>r}\|_K^2
    = \hat\lambda_1 \sum_{j > r} \langle \bv_j, \by\rangle^2/\hat\lambda_j$,
    which is $\hat\lambda_1/\hat\lambda_{r+1}$ times larger than
    $\|\by_{>r}\|_2^2$. Hence $C_1 = \hat\lambda_1/\hat\lambda_{r+1}$
    at worst, which is the \emph{spectral condition number} of the
    head/tail split --- finite only when the head is well-separated
    from the tail.
\item The leverage-incoherence factor $\mu^*$ enters the head term
    via the BLLT-style argument that the empirical Pesenson quality
    $\sigma_\Pi(\Scal_r)$ depends on how spread the eigenvectors are
    over $V$. This is the same machinery as
    \citet{AvronBoutsidis2013} on graph sparsification by leverage
    scores.
\end{enumerate}
Proving Conjecture~\ref{conj:bias} requires the technical tools
listed in \S\ref{sec:tools}, especially (T1)~RKHS norm conversion
and (T3)~leverage-score concentration.
\end{remark}

% =========================================================================
\section{Conjecture: variance bound}
\label{sec:var-conjecture}

\subsection{The principal submatrix sampling problem}

The variance bound requires controlling the random quantity
\begin{equation*}
    \frac{1}{|T^c|}\tr\big(\bK(T^c, T)\bK_T^{+2}\bK(T, T^c)\big)
\end{equation*}
in expectation over $T \sim \Pi_{\mathrm{LSO}}$.

\textbf{The defect in v1 of these notes.}  An earlier draft invoked
BLLT~Lemma~10 (Davis--Kahan + matrix Bernstein) to argue that the
eigenvalues of $\bK_T$ concentrate around scaled population
eigenvalues $\hat\lambda_i \cdot n/N$. This is a \emph{category
error}: BLLT~Lemma~10 is stated for empirical covariance matrices
$\hat\Sigma = \frac{1}{n}\bX^\top\bX$ of i.i.d.\ Gaussian designs
$\bX \in \R^{n \times p}$, where the rows are independent and the
expectation $\E[\hat\Sigma] = \Sigma$ is the population covariance.
The matrix $\bK_T$ is instead a \emph{principal submatrix} of a
fixed Gram matrix $\bK$ extracted by sampling rows/columns without
replacement; it is not an empirical average of independent rank-1
terms, and the population eigenvalues are not its expectation in
any natural sense. The correct theory --- matrix Chernoff for
principal submatrices --- requires \emph{leverage-score} or
\emph{incoherence} conditions on $\bV$, absent from BLLT.

\subsection{What the correct bound looks like}

\begin{conjecture}[Variance bound, transductive analogue of BLLT~\S5.2]
\label{conj:var}
Suppose A\ref{a:noise-bllt}, the leverage incoherence at level
$r^*$ is bounded ($\mu_{r^*}(\bV) \le \mu^*$), and the kernel is
in the BLLT-benign regime: $\reff_{r^*}(\bK)/n \ge c$ and
$\Reff_{r^*}(\bK) \ge \tilde c\, n$ for absolute constants
$c, \tilde c > 0$. Then there exist absolute constants $C_3, C_4$
(depending on the sub-Gaussian norm of $\bxi$, the incoherence
$\mu^*$, and the spectral-decay constants $c, \tilde c$) such that
\begin{equation}
    \E_T\!\big[\mathrm{Var}_{T^c}^2(\sigma_\xi^2, T)\big]
    \;\le\; C_3\,\sigma_\xi^2\frac{r^*}{n}
        \;+\; C_4\,\sigma_\xi^2 \frac{n}{\Reff_{r^*}(\bK)}.
    \label{eq:var-conjecture}
\end{equation}
\end{conjecture}

\begin{remark}[Why this is hard]
\label{rem:var-discussion}
Conjecture~\ref{conj:var} faces three distinct technical
obstacles (cf.\ \S\ref{sec:tools}):
\begin{enumerate}[label=(\arabic*),leftmargin=*]
\item \emph{Spectrum of $\bK_T$ vs.\ spectrum of $\bK$.} The
    eigenvalues of a principal submatrix interlace those of the
    parent matrix (Cauchy interlacing), but the gap depends on
    leverage scores. Quantitative control requires
    matrix-Chernoff-style bounds for sampled
    principal submatrices \citep{Tropp2012,Cohen2016sparsification}.
\item \emph{$T$-dependence of the trace.} The trace in
    \eqref{eq:var-rep} is a function of $T$, not a sum of
    $T$-indexed terms with $T$-independent summands. Standard
    without-replacement concentration (Bardenet--Maillard,
    Hoeffding--Serfling) does not directly apply. The right tool
    is exchangeable-martingale concentration or a stability-based
    argument \`a la \citet{BousquetElisseeff2002}.
\item \emph{Coupling of head and tail.} The BLLT split at level
    $r^*$ assumes the head eigenvectors of $\bK_T$ approximately
    align with the population $\{\bv_j\}_{j \le r^*}$, which holds
    \`a la Davis--Kahan only when the spectral gap
    $\hat\lambda_{r^*} - \hat\lambda_{r^*+1}$ is large relative to
    the perturbation $\|\bK_T - n\bK/N\|_{\mathrm{op}}$. Bounding the
    perturbation requires (1) above.
\end{enumerate}
None of (1)--(3) is fatal; each is a known programme of the
sketching / random-matrix-theory literature
\citep{MahoneyDrineas2009,DerezinskiLiaoDoblerMahoney2020}. But the
adaptation to the leave-node-out setting --- where ``random'' is
the uniform-without-replacement partition $V = T \sqcup T^c$ --- has
not, to our knowledge, been written down.
\end{remark}

% =========================================================================
\section{Conjecture: main result}
\label{sec:main-conjecture}

\begin{conjecture}[Transductive benign-overfitting floor]
\label{conj:main}
Under Conjectures~\ref{conj:bias} and \ref{conj:var}, with
matching hypotheses on $(\by, \bK, T)$, the held-out excess risk of
the min-norm RKHS interpolating learner satisfies
\begin{equation}
    \E_{T,\bxi}\!\big[L_{T^c}(\hat f^{\mathrm{interp}}_\Acal(T), \by)\big]
    \;\le\;
    C_1\,\tau_{r^*}\Var(\by) + C_2\frac{r^*\mu^*}{n-r^*}\Var(\by)
    \;+\; C_3\sigma_\xi^2\frac{r^*}{n}
    \;+\; C_4\sigma_\xi^2\frac{n}{\Reff_{r^*}(\bK)}.
    \label{eq:main-conjecture}
\end{equation}
\end{conjecture}

\begin{remark}[Two-sided window with notes/01 Theorem 1]
\label{rem:window}
Assuming Conjecture~\ref{conj:main} holds, it pairs with
Theorem~1 of \texttt{notes/01} applied to $\Scal_{r^*}$:
\begin{equation*}
    \tau_{r^*}\Var(\by)
    \;\le\; \E[L_{T^c}(\hat f^{\mathrm{interp}}_\Acal(T))]
    \;\le\; C_1\tau_{r^*}\Var(\by) + \text{(slack terms above)}.
\end{equation*}
The window is informative when $C_1 \approx 1$ (i.e., when the
spectral condition number $\hat\lambda_1/\hat\lambda_{r^*+1}$ is
bounded; see Remark~\ref{rem:bias-discussion}) and the slack terms
are small relative to $\tau_{r^*}\Var(\by)$.
\end{remark}

% =========================================================================
\section{The missing tools: a checklist}
\label{sec:tools}

This section catalogues the precise technical ingredients required
to upgrade Conjectures~\ref{conj:bias}--\ref{conj:main} to theorems.

\subsection{(T1) RKHS--$\ell_2(V)$ norm conversion under spectral
truncation}

\begin{problem}
\label{prob:norm-conv}
Establish: for any $f \in \range(\bK)$ orthogonal to $\Scal_r$ in
$\ell_2(V)$, the RKHS norm satisfies
$\|f\|_K^2 \le \hat\lambda_{r+1}^{-1}\|f\|_2^2$ and (more
delicately) the min-norm interpolant $\bb_T(f)$ of $f|_T$ satisfies
\begin{equation*}
    \|\bb_T(f) - f\|_2^2 \;\le\; G(r, T; \bK) \cdot \|f\|_2^2,
\end{equation*}
for an explicit factor $G(r, T; \bK)$ involving the head/tail
spectral gap and the leverage scores.
\end{problem}

The first inequality is a one-line exercise; the second is the
content of BLLT~Lemma~7 in the inductive setting and requires
non-trivial work in the transductive case.

\subsection{(T2) Matrix Chernoff for principal submatrices sampled
without replacement}

\begin{problem}
\label{prob:matrix-chernoff}
Establish: for $\bK \in \R^{N \times N}$ PSD with leverage incoherence
$\mu_r(\bV) \le \mu^*$, and $T \subset V$ uniform of size $n$,
with probability $\ge 1 - \eta$,
\begin{equation*}
    \big\| \bK_T - (n/N) \bK_{T \cup T^c}^{\mathrm{block-diag}} \big\|_{\mathrm{op}}
    \;\le\; C(\mu^*, \eta) \cdot \hat\lambda_1 \cdot \sqrt{\tfrac{n \log N}{N}},
\end{equation*}
or a suitable analogue that gives Davis--Kahan-style alignment
between the top-$r$ eigenvectors of $\bK_T$ and those of $\bK$.
\end{problem}

This is the central missing piece. Relevant literature:
\citet{Tropp2012} for matrix Chernoff/Bernstein,
\citet{Cohen2016sparsification} for leverage-score sampling of
graph Laplacians, \citet{DerezinskiLiaoDoblerMahoney2020} for
spectral analysis of subsampled kernel matrices, and
\citet{BachJordan2002} on Nystr\"om-type approximations. None of
these is directly the statement we need, but each provides a
building block.

\subsection{(T3) Exchangeable-martingale concentration for trace
functionals of $T$}

\begin{problem}
\label{prob:exch-conc}
Establish: for $F : \binom{V}{n} \to \R$ a trace functional of the
form $F(T) = \tr(\bA \bK_T^{+k} \bB)$ with $\bA, \bB$ fixed
matrices and $k \ge 1$, derive a concentration inequality of the
form
\begin{equation*}
    \PR_T[|F(T) - \E_T[F(T)]| \ge t]
    \;\le\; 2\exp\!\big(-t^2/(2 V_F)\big)
\end{equation*}
for a variance proxy $V_F$ tracked in terms of
$\|\bK\|_{\mathrm{op}}$, $\|\bA\|, \|\bB\|$, the spectrum of $\bK$,
and $n, N$.
\end{problem}

This is where Bardenet--Maillard breaks: their inequality concerns
sums of $T$-indexed terms with deterministic summands. For
$T$-dependent summands, the right framework is
\citet{BousquetElisseeff2002}'s algorithmic stability or McDiarmid
on the symmetric group of size-$n$ subsets, via the Aldous--Diaconis
exchangeable-pairs technique \citep{Chatterjee2007exchangeable}.

\subsection{(T4) Tracking BLLT constants in the transductive
setting}

\begin{problem}
\label{prob:bllt-constants}
Adapt the explicit constants of
\citet[][\S5--6]{BartlettLongLugosiTsigler2020} to the
transductive setting, taking into account the
without-replacement corrections from (T2)--(T3).
\end{problem}

% =========================================================================
\section{Application: when can transductive benign overfitting hold?}
\label{sec:application}

Even granting Conjectures~\ref{conj:bias}--\ref{conj:main}, the
BLLT-benign hypothesis (iii) of $\reff_{r^*}/n \to \infty$ is
restrictive. For NTK on common graph families:

\begin{itemize}[label=$\bullet$,leftmargin=*]
\item \emph{Erdős--Rényi or random regular graphs} with edge
    density bounded away from $0, 1$: NTK spectrum has a
    constant-dominated tail above a sharp head; $\reff/n \to \infty$
    and $\mu^* = O(\log N)$. Benign overfitting plausibly holds.
\item \emph{Path / cycle / grid graphs}: NTK spectrum decays
    polynomially $\hat\lambda_j \asymp j^{-\alpha}$; $\reff_r \asymp r$
    and $\Reff_r \asymp r$, so $\reff/n \le 1/2$ for $r \le n/2$.
    Benign overfitting \emph{fails}.
\item \emph{Stochastic block models}: head eigenvectors carry the
    community structure; tail is flat above the spectral gap.
    Benign overfitting holds in the flat-tail regime.
\end{itemize}

For materials informatics (the motivating MatBench application of
the parent paper), the NTK spectrum is typically polynomial
(\`a la \citealp{BordelonCanataPehlevan2020}), placing it in the
non-benign regime. The transductive benign-overfitting floor of
Conjecture~\ref{conj:main} would not directly apply; the
operative bound there is instead Theorem~1 of \texttt{notes/01}
with explicit regularisation
($\Omega = \|\cdot\|_K + \lambda\|\cdot\|_2$) as in the ridge
analysis of \citet{TsiglerBartlett2020}. Pursuing this ridged
variant is a separate research direction.

% =========================================================================
\section{Status and references}
\label{sec:status}

\subsection{What this note is and is not}

\begin{itemize}[label=$\bullet$,leftmargin=*]
\item It \textbf{is} a precise statement of the
    transductive-benign-overfitting research programme, with the
    bias--variance decomposition proven (\S\ref{sec:decomp}), the
    target bound conjectured (\S\S\ref{sec:bias-conjecture}--\ref{sec:main-conjecture}),
    and the missing technical tools enumerated (\S\ref{sec:tools}).
\item It \textbf{is not} a proof of any new theorem. The
    Conjectures of \S\S\ref{sec:bias-conjecture}--\ref{sec:main-conjecture}
    should not be cited as established results.
\item It \textbf{supersedes} an earlier draft
    (\texttt{notes/08\_benign\_overfitting\_NTK.tex}, v1) which
    contained defective ``proof sketches'' for the same
    Conjectures, identified by adversarial peer review. The
    defects --- norm-conversion confusion, category error on
    BLLT~Lemma~10, conditional-expectation error on Bardenet--Maillard
    --- are documented in the header comments of this file.
\end{itemize}

\subsection{Implication for \texttt{notes/07}}

The forward pointer in \texttt{notes/07}~\S5.3 (Axis~4 operational
difficulty) should reference this note as a \emph{research
programme}, not as a result. Specifically, the sentence
``\emph{the operational NTK/Mercer-truncated bound is non-trivial
and deferred to notes/08...}'' should be qualified as
``\emph{deferred to notes/08, which formulates this as a precise
sequence of conjectures parameterised by missing technical tools}.''

\subsection{Implication for \texttt{paper.tex}}

\texttt{paper.tex}~\S8.3 item~(4) should be updated to acknowledge
that, for over-parameterised NN / NTK learners with strictly PD
kernels, the universal spectral floor is operationally trivial
(reduces to $0$ when $\range(\bK) = \R^N$), and that a non-trivial
floor for this setting is an open problem rather than a result of
the framework. The BLLT-style adaptation is a natural attack but
the technical obstacles enumerated in this note must be resolved
first.

% =========================================================================
\bibliography{../references/references}

\end{document}
